\documentclass[10pt]{article}

\usepackage[utf8]{inputenc}

\usepackage[T1]{fontenc}

\usepackage{amsmath}

\usepackage{amsfonts}

\usepackage{amssymb}

\usepackage[version=4]{mhchem}

\usepackage{stmaryrd}

\usepackage{bbold}



\begin{document}

где $\bar{S}_{n}=\max \left(0, S_{1}, \ldots, S_{n}\right)$. Факторизационные тождества представляют собой мощное средство изучения граничных задач для С. б. (см. [1], т. 2; [2], [3]).



В настоящее время граничные задачи для С. б. исследованы весьма полно, включая асимптотический их анализ (см. [1], т. 2; [4]-[6]).



Аналитически решение граничных задач приводит к интегрально-разностным уравнениям. Напр., вероятность $u_{n}(x, a, b)$ того, что частица, вышедшая из точки $x \in(-b, a)$, покинет за время $n$ интервал $(-b, a)$, удовлетворяет уравнению (формула полной вероятности по первому скачку)



\[

\begin{gathered}

u_{n+1}(x, a, b)=\int_{-F}^{a-x}(a-x)+F(-b-x)

\end{gathered}

\]



где $F(x)=\mathrm{P}\left\{X_{1}<x\right\}$. Если перейти к производящим функциям $u(z, x, a, b)=\sum_{n=1}^{\infty} z^{n} u_{n}(x, a, b)$, то получают обычные ивтегральные уравнения. Существует, по крайней мере, два подхода к исследованию асимптотич. свойств решений этих уравнений. Один из них основан на изучении аналитич. свойств двойного преобразования



\[

U(z, \lambda, a, b)=\int e^{i \lambda x} u(z, x, a, b) d x

\]



и последующего его обращевия (см. [4] - [6]). Другой связан с привлечением методов Вишика и Люстерника решения уравнений с малым параметром (см. [6]). На втором пути обнаруживаются глубокие связи рассматриваемых задач с теорией потенциала.



Многое из сказанного выше переносится в на С. б. с зависимыми скачками, когда случайные величины $S_{n}$ связаны в цепь Маркова, а также на многомерные C. б. в $\mathbb{R}^{d}, d>1$ (см. [6], [7]).



Jium.: [1] $\Phi$ е л л е p В., Введение в теорию вероятностей и ее приложения, пер. с англ., 2 изд., т. 1, М., 1967; [2] Б о р о вк о в А. А., Теория вероятностей, М., 1976; [3] е г о ж е, Вероятностные процессы в теории массового обслуживания, М., 1972 ; [4] К ор о л ю к В. С., Б о р о в с к и х Ю. В., Аналитические проблемы асимштотики веронтностных распределений, К., 1981; [5] Б ор о в к о в А. А., «Сиб. матем. ж.», 1962, т. 3, № 5, с. 645-94; [6] JI о т о в В. И., «Теория вероятн. И ее примен.», 1979 , т. 24 , № 3, с. $475-85 ;$ № 4, с. $873-79 ;$ [7] С II пц е р Ф., Принципы случайного блуждания, пер. с англ., М.,

1969 .



СЈУ ЧАИНОЕ КОДИРОВАНИЕ - один из методов кодирования (см. Кодирование и декодирование), при к-ром каждому возможному значению сообщения, вырабатываемому источником сообцений, ставится в соответствие случайно выбранное значение сигнала на входе канала связи. При әтом на множестве значений сигналов на входе канала задается нек-рое распределение вероятностей. Часто предполагается, что каждый элемент кода (т. е. значение сигнала на входе, соответствующее данному значению сообщения) выбирается независимо от других и в соответствии с данным распределением вероятностей. Иногда С. к. определяют так, чтобы каждая реализация С. к. была групповЫМ Кодом.



Важность рассмотрения С. к. связана с тем обстоятельством, что осредненная по всем реализациям очибочного декодирования вероятность дает относительно легко исследуемую оценку сверху для вероятности ошибочного декодирования оптимального кода.



Лum.: [1] ШI е н н о н К., Работы по теории информации и кибернетике, пер. с англ., М. 1963 , с. $243-332$; [2] Д о б р уш п н Р. Л., "Успехи матем. наук», 1959, т. 14, в. 6, с. $3-104 ;$ пер. с англ., М., 1974. Р. Л. Добрушия, В. В. Прелов.



СЛУЧАИНОЕ ОТОБРАЖЕНИЕ $\sigma$ множества $X=$ $=\{1,2, \ldots, n\}$ в себя-случайная величина, принимающая значения из множества $\sum_{n}$ всех однозначных отображений множества $X$ в себя. С. о. $\sigma$, Для к-рых вероятность $\mathrm{P}\{\sigma=s\}>0$ только для взаимно однозначных отображений $s \in \sum_{n}$, наз. с л у ч а й н ы м и п о д с т а н о в к а м и степени $n$. Наиболее полно изучены С. о., для к-рых $\mathrm{P}\{\sigma=s\}=n^{-n} п \mathrm{pu}$ всех $s \in \sum_{n}$. Реализация такого С. о. представляет собой результат простого случайного выбора из множества $\sum_{n}$.



Лuт.: [1] К о л ч в н В. Ф., Случайные отображения, М., CІУчАЙО ПОЛЕ, С. Колчин. мн о гоме рным в р е е нем, или с многом е р ны м п а р а м т р о м,- случайная функция, заданная на множестве точек какого-то многомерного пространства. С. П. представляют собой важный тип случайных бункций, часто встречающийся в различных приложениях. Примерами С. п., зависящих от трех пространственных координат $x, y, z$ (а также и от времени $t$ ), могут служить, в частности, поля компонент скорости, давления и температуры турбулентного төчения жидкости или газа (см. [1]); С. І., зависящим от двух координат $x$ и $y$, будет высота $z$ взволнованной морской поверхности или поверхности какой-либо пероховатой пластинки (см. [2]); при исследовании глобальных атмосферных процессов в масштабе всей Земли поля наземного давления и др. метеорологич. характеристик иногда рассматриваются как С. II. на сфере и т. д.



Теория С. п. общего вида фактически не отличается от общей теории случайных функций; более содержательные конкретные результаты удается получить лишь для ряда специальных классов С. п., обладающих дополнительными свойствами, облегчающими их изучение. Одним из таких классов является класс случайных полей однородных, заданных на однородном пространстве $S$ с группой преобразований $G$ и обладающих тем свойством, что распределения вероятностей значений поля на произвольной конечной группе точек пространства $S$ или же среднее значение поля и вторые моменты его значений в парах точек не меняются при применении к аргументам поля какого-либо преобразования из групшы $G$. Однородные С. п. на евклидовом пространстве $\mathbb{R}^{k}, \quad k=1,2, \ldots$, или на решетке $\mathbb{Z}^{k}$ точек $\mathbb{R} k$ с целочисленными координатами, отвечающие выбору в качестве групшы $G$ совокупности всевозможных (или всех целочисленных) параллельных переносов, являются естественным обобщением стачионарных случайных процессов, на к-рое просто переносится бо́льшая часть результатов, доказанных для таких процессов; большой интерес для приложений (в частности, для механики турбулентности, ср. [1]) представляют также т. н. однородные и изотропные поля на $\mathbb{R}^{3}$ и $\mathbb{R}^{2}$, отвечающие выбору в качестве группы $G$ совокупности всевозможных изометрич. преобразований соответствующего пространства. Важной особенностью однородных С. П. является существование спектральных разложений специального вида как самих таких полей, так и их корреляционных функций (cм., напр., [3], [4]).



Другим привлекающим много внимания классом С. п. является класс м а р к о в с ки х с л уч а й ны п о л е й, заданных в нек-рой области $K$ пространства $\mathbb{R}^{k}$. Условие марковости С. ІІ. $U(\boldsymbol{x})$, грубо говоря, означает, что для достаточно широкой совокупности открытых множеств $Q$, имеющих границу $\Gamma$, фиксация значений поля в $\varepsilon$-окрестности $\Gamma^{\varepsilon}$ границы $\Gamma$ при любом $\varepsilon>0$ делает семейства случайных величин $\left\{U(x), x \in Q \backslash \Gamma^{\varepsilon}\right\}$ и $\left\{U(x), x \in T \backslash \Gamma^{\varepsilon}\right\}$, где $T-$ дополнение замыкания $Q$ в $K$, взаимно независимыми (или, в случае марковости в широком смысле, взаимно некоррелированными; см., напр., [5]). Обобщением понятия марковского С. П. является понятие $L$-марковского С. П., для к-рого указанная выше независимость (или некоррелированность) имеет место лишь при замене границы $\Gamma$ области $Q$ специальным образом оп-





\end{document}